\documentclass[11pt,a4paper]{article}

\usepackage[margin=1in]{geometry}
\usepackage{amsmath,amssymb,amsthm}
\usepackage{mathrsfs}
\usepackage{booktabs}
\usepackage{microtype}
\usepackage[hidelinks]{hyperref}

\newtheorem{remark}{Remark}

\newcommand{\PP}{\mathbb{P}}
\newcommand{\C}{\mathbb{C}}
\newcommand{\R}{\mathbb{R}}
\newcommand{\Z}{\mathbb{Z}}
\newcommand{\Vol}{\operatorname{Vol}}
\newcommand{\Sym}{\operatorname{Sym}}
\newcommand{\Stab}{\operatorname{Stab}}
\newcommand{\rms}{\operatorname{rms}}
\newcommand{\dd}{\,\mathrm{d}}
\newcommand{\FS}{\mathrm{FS}}

\title{Approximate Calabi--Yau metrics on the Fermat quartic:\\
       method and current residuals}
\author{}
\date{\today}

\begin{document}
\maketitle

\begin{abstract}
We report the state of the numerical approximation of the Ricci-flat metric on
the Fermat quartic $X=\{Z_0^4+Z_1^4+Z_2^4+Z_3^4=0\}\subset\PP^3$, in the
algebraic ansatz of Headrick--Nassar. The best converged residual is
$\operatorname{mean}|\eta-1|=1.22\cdot10^{-12}$ at degree $k=24$ and
$1.9\cdot10^{-13}$ at $k=28$. The observed convergence rate is $2.68$ to $2.81$
per unit of $k$ over $k=12$ to $24$, against the value $2.83$ that the
Headrick--Nassar energy law predicts.
\end{abstract}

\section{The ansatz and its symmetry}

Let $\{s_A\}$ be the monomial basis of $H^0(X,\mathcal{O}(k))$ and put
\begin{equation}
  \varphi=\frac1k\log p,
  \qquad
  p=\sum_{A,B}h_{A\bar B}\,s_A\bar s_B,
  \qquad
  \omega=i\partial\bar\partial\varphi ,
  \label{eq:ansatz}
\end{equation}
with $h$ a positive Hermitian matrix. Write
$\mu_\omega=\tfrac1{2}\omega^2$ for the volume form of $\omega$,
$\mu_\Omega=\Omega\wedge\bar\Omega$ for that of the holomorphic volume form, and
\begin{equation}
  \hat\mu=\frac{\Vol(X,\mu_\omega)}{\Vol(X,\mu_\Omega)}\,\mu_\Omega,
  \qquad
  \eta=\frac{\mu_\omega}{\hat\mu},
  \qquad
  f=1-\frac1\eta .
  \label{eq:eta}
\end{equation}
The Calabi--Yau equation is $\eta\equiv1$, so $f$ measures the failure of
$\omega$ to be Ricci-flat.

The symmetry group of $X$ is
$\Gamma=(\Z_4^3\rtimes S_4)\times\Z_2$, where the last factor is complex
conjugation; the holomorphic part has order $1536$. We impose $\Gamma$-invariance
on $h$ exactly, and this is what makes the fit small.

\paragraph{Block structure.}
Invariance under the phase group $\Z_4^3$ forces $h_{A\bar B}=0$ unless
$A\equiv B \pmod 4$. Hence $h$ is block diagonal in the residue $r=A\bmod4$, with
$64$ blocks. Inside a block $A=r+4C$, so with
\begin{equation}
  u_a:=z_a^4,
  \qquad
  z^A=z^r u^C ,
\end{equation}
the sections of the block are the degree-$m_r$ monomials $u^C$, where
$m_r=(k-|r|)/4$, and the defining equation of $X$ becomes exactly
\begin{equation}
  \textstyle\sum_a u_a=0 .
  \label{eq:relation}
\end{equation}
Equation \eqref{eq:relation} is the one linear relation among the block
monomials, and it is the complete effect of the fact that the $s_A$ are not
independent on $X$. The block section space is therefore
$\Sym^{m}(\mathrm{std})$, with $\mathrm{std}$ the $3$-dimensional standard
representation of $S_4$, of dimension $d_r=\binom{m_r+2}{2}$. As a check,
\begin{equation}
  \sum_r \binom{m_r+2}{2}=2k^2+2
  \label{eq:seccount}
\end{equation}
holds identically; we verified it at $k=4,8,12,32,38$.

\paragraph{Parameter count.}
$S_4$ permutes the blocks, so one parameter set belongs to each $S_4$ orbit of
residues, constrained by the stabiliser, and conjugation makes $h$ real
symmetric. Therefore
\begin{equation}
  n_\theta=\sum_{\text{orbits}}\dim\Sym^2\!\bigl(\Sym^{m_r}(\mathrm{std})\bigr)^{\Stab(r)} ,
  \label{eq:ntheta}
\end{equation}
which we evaluate exactly by character theory:

\begin{center}
\begin{tabular}{lrrrrrrr}
\toprule
$k$ & 12 & 16 & 20 & 24 & 28 & 32 & 38\\
\midrule
sections $2k^2+2$ & 290 & 514 & 802 & 1{,}154 & 1{,}570 & 2{,}050 & 2{,}890\\
$n_\theta$        &  56 & 135 & 285 &   542 &   949 & 1{,}559 & 2{,}987\\
\bottomrule
\end{tabular}
\end{center}

\noindent
So the invariance reduces $2k^2+2$ sections to $n_\theta\approx k^2/0.85$ real
parameters; at $k=28$ that is $949$ against the $1{,}570^2$ entries of an
unconstrained Hermitian $h$.

\section{Numerical method}

\paragraph{Points.}
We sample $X$ by intersecting it with random projective lines, each line giving
four points, and assign the mass $\mu_\Omega/\mathrm{vol}_\FS$. Two Newton steps
refine each point onto $X$: over the sets used below,
$\max|P(z)|/\sum_a|z_a|^4=3.3\cdot10^{-16}$ with median $7.3\cdot10^{-17}$.
The fit set and the reported set are independent samples with different seeds.

\paragraph{Evaluation of the metric.}
We restrict to $X$ by sub-patches chosen by the largest gradient of $P$, with a
per-point restriction matrix computed once, and then
$v=\det g|_X/\mu_\Omega$ and $\eta=v/\langle v\rangle$. Three choices matter
below $10^{-13}$ and each was measured:
\begin{itemize}
  \item the monomial exponents are kept as integers, so $z^A$ never passes
        through $\exp$ and $\log$;
  \item the multinomial coefficients are formed in exact integer arithmetic,
        since $k!/\prod_a A_a!$ overflows \texttt{int64} from $k=24$;
  \item $\varphi_{i\bar j}$ is evaluated in the centred Gram form
        $t_i=\partial_i s-(\langle\partial_i s,s\rangle/p)\,s$,
        $\varphi_{i\bar j}=\langle t_i,t_j\rangle/p$, which avoids the
        cancellation of the difference form $(p_{i\bar j}p-p_i\bar p_j)/p^2$.
\end{itemize}
Independently, $v$ agrees with the bordered-determinant formula of
Headrick--Nassar, and it agrees between the two sub-patch choices at points where
two gradient components are nearly equal, in both cases to $10^{-15}$.

\paragraph{Loss and solver.}
We minimise the mass-weighted mean square percentage error
\begin{equation}
  E=\sum_i w_i\Bigl(\frac{\eta_i}{\langle\eta\rangle}-1\Bigr)^2,
  \qquad
  w_i=\frac{\text{mass}_i}{\sum_j \text{mass}_j},
  \label{eq:loss}
\end{equation}
with the Levenberg--Marquardt implementation of \texttt{optimistix}. $E$ is one
of the energy functionals of Headrick--Nassar, convex with its minimum at
$\eta=1$.

The Jacobian is analytic rather than nested autodiff. Since $p$ is linear in
$\theta$ and $v$ is built from Gram sums of the section vector and its four
holomorphic derivatives, Jacobi's formula plus the contraction with
$K=R^\dagger g^{-1}R$ gives $\partial v/\partial\theta_\alpha$ as one Frobenius
product per block. This is $126$ times faster than \texttt{jacfwd} at $k=20$.
The residual is evaluated in chunks over the points, so the peak memory follows
the chunk and not the size of the fit set.

\paragraph{Conditioning and precision.}
Two further choices are needed at large $k$:
\begin{itemize}
  \item $\theta$ is divided by a single global factor, so the fit starts at
        $\pm1$. The invariant basis is orthonormal in the Frobenius norm, so the
        raw components of $\theta_\FS$ reach $10^{14}$ at $k=28$, and
        Levenberg--Marquardt damps and tests steps in absolute size.
  \item the Gram sums $p$, $\partial p$ and $\partial\bar\partial p$ are
        accumulated in double-double arithmetic. Against a reference in
        \texttt{mpmath} at $50$ digits this lowers the relative error of
        $\varphi$ at $k=28$ from $9.9\cdot10^{-14}$ to $3.2\cdot10^{-15}$, and
        the corresponding floor on $|f|$ from $4.0\cdot10^{-13}$ to
        $1.6\cdot10^{-14}$. At the fitted $\theta$ the measured relative error of
        $v$ is $4.4\cdot10^{-16}$ at $k=24$ and $2.0\cdot10^{-15}$ at $k=28$.
\end{itemize}
The bases are made canonical by Gram--Schmidt in a fixed order, rather than by a
singular value decomposition of a degenerate subspace, and each saved $\theta$
carries a fingerprint of the convention it belongs to.

Optionally the section basis is whitened in $L^2(X)$: with $G$ the Gram matrix of
the block sections on $X$ and $W=G^{-1/2}$, one parametrises $h=W h_w W$. This is
an exact reparametrisation, and one matrix per orbit suffices because
$G$ for any block is conjugate to that of its orbit representative by the same
rotation that relates the blocks. It reduces the condition number of the
Jacobian by $576$ times at $k=24$ and leaves the residual unchanged; we use it
when the coefficients themselves are reported.

\section{Results}

All residuals below are computed on a held-out set of $8\cdot10^6$ points, for
which the standard error of $\operatorname{mean}|f|$ is $0.04\%$. The
normalisation of $\eta$ is the sampled mass-weighted mean of $v$. Norms use
$\Vol(X,\mu_\omega)=2\pi^2$, and
$\Vol(X,\mu_\Omega)=\Gamma(1/4)^8/(64\pi^2)=47.268180032308\ldots$

\begin{center}
\begin{tabular}{rrllllrr}
\toprule
$k$ & $n_\theta$ & $\operatorname{mean}|f|$ & $\rms(f)$
    & $\max|f|$ & $\|f\|_{L^2}$ & fit pts & rate\\
\midrule
12 &  56 & $2.202\cdot10^{-7}$  & $3.732\cdot10^{-7}$  & $1.116\cdot10^{-5}$  & $1.658\cdot10^{-6}$  & 100{,}000 & ---\\
16 & 135 & $3.536\cdot10^{-9}$  & $6.701\cdot10^{-9}$  & $2.865\cdot10^{-7}$  & $2.977\cdot10^{-8}$  &  27{,}000 & 2.81\\
20 & 285 & $6.323\cdot10^{-11}$ & $1.332\cdot10^{-10}$ & $7.308\cdot10^{-9}$  & $5.918\cdot10^{-10}$ &  57{,}000 & 2.73\\
24 & 542 & $1.220\cdot10^{-12}$ & $2.800\cdot10^{-12}$ & $1.741\cdot10^{-10}$ & $1.244\cdot10^{-11}$ & 108{,}400 & 2.68\\
28 & 949 & $1.9\cdot10^{-13}$   & $2.6\cdot10^{-13}$   & $8.2\cdot10^{-12}$   & $1.1\cdot10^{-12}$   &  59{,}312 & 1.6\\
\bottomrule
\end{tabular}
\end{center}

\noindent
The $k=28$ row is the median of six fits; see Remark~\ref{rem:k28}. Every fit
converged, and the metric stays well inside the positive cone: at $k=28$ the
smallest eigenvalue of $\varphi_{i\bar j}$ over the held-out set is $3.496$ and
$\min p=6.6\cdot10^{9}$.

\paragraph{Convergence rate.}
Over $k=12$ to $24$ the residual falls by a factor $2.68$ to $2.81$ per unit of
$k$. The energy law of Headrick--Nassar for the quartic, $E\approx0.03\cdot8^{-k}$
with $E\sim\sigma^2$, predicts $2.83$. Fitting
$\operatorname{mean}|f|\approx A\,\lambda^{-k}$ to $k=12$ and $16$ gives
$\lambda=2.80$ and $A=5.2\cdot10^{-2}$, and this reproduces $k=20$ and $k=24$ to
$6\%$ and $26\%$.

\paragraph{Shape of the residual.}
The ratios $\rms(f)/\operatorname{mean}|f|$ and $\max|f|/\rms(f)$ are
$1.89,\,2.11,\,2.30$ and $42.8,\,54.9,\,62.2$ at $k=16,20,24$. Both grow slowly
with $k$, so the residual becomes less uniform as it becomes smaller.
We stress that $\max|f|$ is the largest value over a finite sample and is
therefore a lower bound on $\sup_X|f|$; it grows with the size of the sample and
must not be read as a bound.

\begin{remark}[the residual at $k=28$]
\label{rem:k28}
At $k=28$ the value reached is not reproducible to better than a factor of
$2.6$: six fits from the same starting point, on the same $59{,}312$ fit points,
differing only in the order of summation over the points, span
$1.12\cdot10^{-13}$ to $2.93\cdot10^{-13}$, with median $1.91\cdot10^{-13}$. The
held-out residual tracks the fit-set residual to $1\%$ in every case, so this is
not overfitting, and the fits use their entire step budget, so it is not early
termination. The condition number of the Jacobian is $6\cdot10^{17}$ at $k=28$,
so the minimum lies in a valley that is flat over many directions, and the
iteration halts where the computed gradient falls to round-off, which occurs at a
range of heights rather than at one. We therefore quote the median and the range,
and we regard the rate $1.6$ for $24\to28$ as a property of the arithmetic and
not of the ansatz.
\end{remark}

\begin{remark}[precision of the coefficients]
The residual is a forward evaluation on held-out points and does not involve the
Jacobian; it is unchanged to five significant figures when the condition number
of the Jacobian is altered by $576$ times. The coefficients $h$ are a different
matter. With $\kappa$ the condition number of the Jacobian, the relative error of
$\theta$ is of order $10^{-18}\kappa$, so about four digits of $h$ survive at
$k=24$ in the whitened parametrisation, and essentially none at $k=28$. Reaching
the residual is thus much easier than determining the metric that attains it;
the latter needs extended precision in the linear solve.
\end{remark}

\section{Outlook}

The residual floor imposed by the arithmetic is the binding constraint above
$k\approx28$. With the Gram sums alone in double-double, the measured floor on
$|f|$ grows as $1.28$ per unit of $k$ while the residual falls as $2.68$, so the
margin closes near $k=29$. Extending the compensated accumulation to the
remaining stages of the evaluation flattens the floor in $k$ at approximately
$1.6\cdot10^{-15}$, which would place $\operatorname{mean}|f|\sim10^{-15}$ within
reach at $k\approx31$.

Two quantities remain to be measured and are needed for the perturbation
argument: the ratio $\|\nabla f\|_{L^2}/\|f\|_{L^2}$, which converts the
residuals above into the $H^1$ norm in which the requirement is stated, and a
rigorous enclosure of $\sup_X|f|$ in place of the sampled maximum.

\end{document}
